\documentclass[../../../include/open-logic-section]{subfiles}

\begin{document}

\olfileid{sth}{choice}{countablechoice}
\olsection{Pilihan Kaétung}

Tanpa migunakaké Pilihan utawa Pengurutan Becik, kita gampang
mbuktekaké pratelan iki:

\begin{lem}[ing $\Zminus$]
Saben himpunan winates nduwèni fungsi pilihan.
\end{lem}

\begin{proof}
Anggepen $a = \{b_1, \ldots, b_n\}$. Supaya prasaja, anggepen saben
$b_i \neq \emptyset$. Mula ana obyèk $c_1, \ldots, c_n$ kang nyukupi
$c_1 \in b_1, \ldots, c_n \in b_n$. Miturut
\olref[z][pairs]{prop:pairsconsequences}, himpunan $\{\langle b_1,
c_1\rangle , \ldots, \langle b_n, c_n\rangle\}$ ana; iki fungsi
pilihan kanggo~$a$.
\end{proof}

Nanging prakarané dadi luwih ruwet nalika kita nggatekaké himpunan
tanpa wates. Upamané, gatekna perluasan ``paling cilik'' saka asil mau:

\begin{defish}
\emph{Pilihan Kaétung.} Saben himpunan \emph{kaétung} nduwèni fungsi pilihan.
\end{defish}

Iki kasus khusus saka Aksioma Pilihan. Pranyata prinsip iki wis
bola-bali digunakaké tanpa kesadaran kang cetha. Ing ngisor iki ana
rong conto kang migunani.\footnote{Conto iki asalé saka
\citet[\S9.4]{Potter2004} lan Luca Incurvati.}

\begin{ex}
Pamikiran kang lumrah yaiku: kanggo sembarang himpunan $A$, salah siji
$\cardle{\omega}{A}$, utawa $\cardeq{A}{n}$ kanggo sawijining
$n \in \omega$. Iki sawijining cara kanggo mratelakaké pamikiran
intuitif yèn saben himpunan winates utawa tanpa wates. Cantor lan
akèh matématikawan liyané nyatakaké iki tanpa bukti. Amarga ngati-ati,
kita wis mbuktekaké ing
\olref[cardinals][classing]{generalinfinitycharacter}. Nanging bukti
ing kana lumaku ing $\ZFC$: kita nganggep saben himpunan $A$ bisa
diurutaké kanthi becik, mula $\card{A}$ mesthi ana. Dadi kita kanthi
gamblang nganggep Pilihan.

Malah \citet{Dedekind1888} mènèhi buktiné dhéwé kanggo
pratelan iki:

\begin{thm}[ing $\Zminus + \text{Pilihan Kaétung}$]
Kanggo sembarang $A$, salah siji $\cardle{\omega}{A}$ utawa
$\cardeq{A}{n}$ kanggo sawijining $n \in \omega$.
\end{thm}

\begin{proof}
Anggepen $\cardneq{A}{n}$ kanggo saben $n \in \omega$. Mula kanggo
saben $n < \omega$ ana himpunan pérangan $A_n \subseteq A$ kang
nduwèni pas $\cardexpo{2}{n}$ unsur. Saka runtutan $A_0, A_1, A_2,
\ldots$ iki, kanggo saben $n$ kita dhéfinisikaké:
\[
	B_n = A_n \setminus \bigcup_{i < n} A_i.
\]
Saiki gatekna
\begin{align*}
	\card{\bigcup_{i < n}A_i}
	&\leq \card{A_0} + \card{A_1} + \ldots + \card{A_{n-1}}\\
	&=1 + 2 + \ldots + 2^{n-1}\\
	& = 2^n - 1\\
	& < 2^n = \card{A_n}
\end{align*}
Mula saben $B_n$ ora kosong; jupuken sawijining unsur $c_n$ saka
$B_n$. Himpunan-himpunan mau uga pisah-pisah sepasang-sepasang;
mula yèn $c_n = c_m$, banjur $n = m$. Saben $c_n \in A$, saéngga
fungsi $f(n) = c_n$ iku injeksi $\omega \to A$.
\end{proof}
\noindent
Dedekind ora nyebut yèn dhèwèké wis migunakaké Pilihan Kaétung.
Nanging apa \emph{panjenengan} ndeleng panganggoné? Gatekna maneh.
(Pancen, \emph{gatekna maneh}.)

Bukti iki migunakaké Pilihan Kaétung kaping pindho. Kapisan kanggo
milih runtutan himpunan $A_0$, $A_1$, $A_2$, \dots\@; kapindho
kanggo milih unsur $c_n$ saka saben~$B_n$. Panganggoné Pilihan iki
ora bisa disingkiraké. \citet[kaca~138]{Cohen1966} mbuktekaké yèn
asile gagal tanpa sembarang wujud Pilihan. Tegesé, konsisten karo
$\ZF$ yèn ana himpunan kang \emph{ora bisa dibandhingaké} karo~$\omega$.
\end{ex}

\begin{ex}
Ing \citeyear{Cantor1878}, Cantor nyatakaké yèn gabungan kaétung
saka himpunan kaétung uga kaétung. Dhèwèké ora mènèhi bukti,
mbokmenawa amarga nganggep buktiné cetha. Kanthi ngati-ati, kita wis
mbuktekaké versi kang luwih umum ing
\olref[card-arithmetic][simp]{kappaunionkappasize}. Nanging bukti mau
kanthi gamblang nganggep Pilihan. Malah versi kang kurang umum iki
uga mbutuhaké Pilihan Kaétung.

\begin{thm}[ing $\Zminus + \text{Pilihan Kaétung}$]
Yèn $A_n$ kaétung kanggo saben $n \in \omega$, mula $\bigcup_{n <
\omega} A_n$ uga kaétung.
\end{thm}

\begin{proof}
Tanpa ngurangi umumé, anggepen saben $A_n \neq \emptyset$.
Mula kanggo saben $n \in \omega$ ana !!a{surjection}
$f_n \colon \omega \to A_n$. Dhéfinisikaké
$f \colon \omega \times \omega \to \bigcup_{n < \omega} A_n$
kanthi $f(m, n) = f_n(m)$. Asilé katut amarga $\omega \times \omega$
kaétung (\olref[sfr][siz][zigzag]{natsquaredenumerable}) lan $f$
iku !!a{surjection}.
\end{proof}
\noindent
Apa panjenengan ndeleng panganggoné Pilihan Kaétung? Prinsip iki
digunakaké kanggo milih runtutan fungsi $f_0$, $f_1$, $f_2$,
\dots\footnote{Panganggoné Pilihan kang padha uga ana ing
\olref[card-arithmetic][simp]{kappaunionkappasize}, nalika kita
ndhawuhaké ``Kanggo saben $\beta \in \cardfont{a}$, tetepna
!!a{injection} $f_\beta$''.} Tanpa prinsip Pilihan, asil iki uga
bisa gagal. Mliginé, \citet{FefermanLevy1963} mbuktekaké yèn
konsisten karo $\ZF$ menawa himpunan bilangan riil $\Real$, kang
ora kaétung, dadi gabungan kaétung saka himpunan-himpunan kaétung.
Nanging Russell mènèhi gambaran kang luwih ngélingaké:
\begin{quote}
  Bab iki digambaraké déning sawijining milyuner kang saben tuku
  sapasang sepatu bot uga tuku sapasang kaos kaki, lan ora naté
  tuku kaos kaki ing wektu liya. Dhèwèké seneng banget tuku
  kaloroné nganti pungkasané nduwèni $\aleph_0$ pasang bot lan
  $\aleph_0$ pasang kaos kaki\dots\@ Ing saben pasangan bot, kita
  bisa mbédakaké sisih tengen lan kiwa, mula bisa milih siji saka
  saben pasangan: kabèh bot tengen utawa kabèh bot kiwa. Nanging
  ing pasangan kaos kaki ora ana patokan pilihan kang padha; tanpa
  nganggep aksioma ping-pingan [yaiku, kanthi efektif, Pilihan],
  kita ora bisa mesthèkaké anané kelas kang ngemot siji kaos kaki
  saka saben pasangan.
  \citep[kaca~126]{Russell1919}
\end{quote}
Cekaké, kanggo mbuktekaké yèn saka pasangan kaos kaki kang kaétung
cacahé bisa kawangun himpunan kaos kaki kang kaétung, dibutuhaké
sawijining wujud Pilihan. Pancen, tanpa Pilihan Kaétung utawa
prinsip kang ekuivalen, gabungan kaétung saka himpunan kaétung
bisa dadi ora kaétung.
\end{ex}

Piwulangé: Pilihan Kaétung bola-bali digunakaké tanpa akèh
kesadaran saka panganggoné. Pitakonan filsafaté yaiku: kepriyé kita
bisa \emph{mbeneraké} Pilihan Kaétung?

Salah siji upaya pambenaran intuitif bisa njaluk bantuan saka
tugas tanpa wates kang rampung sajroning wektu winates.
Anggepen pilihan kapisan rampung sajroning $\nicefrac{1}{2}$ menit,
pilihan kapindho sajroning $\nicefrac{1}{4}$ menit, \dots, lan
pilihan kaping $n$ sajroning $\nicefrac{1}{2^n}$ menit, \dots\@
Mula sajroning $1$~menit kita wis nindakaké runtutan $\omega$
pilihan lan nemtokaké sawijining fungsi pilihan.

Nanging apa sejatine kang bisa diwulangaké déning eksperimen pikiran
kaya mangkono? Sepisan, eksperimen iku mangertèni ``milih'' kanthi
teges kang harfiah banget. Kajaba iku, matématika kaya dikaitaké
marang kamungkinan metafisis.

Kang luwih wigati, panalaran mau ora bakal mbeneraké Pilihan
\emph{kanthi umum}, nanging mung Pilihan Kaétung \emph{waé}.
Yèn \emph{saben} himpunan kudu nduwèni fungsi pilihan, kita kudu bisa
nindakaké ``tugas tanpa wates kang dawané sembarang ordinal''.
Terus terang, gagasan iku ora masuk akal.

\end{document}
